\hypertarget{appendix-f11d1a521253}{}\label{referencestart:appendix-f11d1a521253}
\section{Floating-point environment and compiler modes}\label{appendix-f11d1a521253-floating-point-environment-and-compiler-modes}

\textbf{Curriculum depth:} Developed · \textbf{First seen:} C004 · \textbf{Latest development:} C005

\begin{allmightroledefinition}

\subsection{Definition and mechanism}\label{appendix-f11d1a521253-definition-and-mechanism}

The environment includes rounding direction and exception status. Runtime mode changes can fail and must restore prior state. Strict compiler settings are needed; fast-math may discard required special-value or environment assumptions.

\end{allmightroledefinition}

M005 uses strict arithmetic and disables implicit FMA contraction to distinguish fused from separate graphs. The full floating environment is restored after the lab; cross-platform bitwise reproducibility is not established.

\subsection{Prerequisites}\label{appendix-f11d1a521253-prerequisites}

\begin{itemize}
\tightlist
\item
  \hyperref[appendix-637d164ae5ca]{Directed rounding modes} --- The signed rounding direction is state that a demonstration or enclosure algorithm must control.
\item
  \hyperref[appendix-d6607cc4c235]{NaNs and payloads} --- Unordered comparisons and payload distinctions require explicit metric, comparison and serialization policy.
\end{itemize}

\subsection{Used by}\label{appendix-f11d1a521253-used-by}

\begin{itemize}
\item
  \hyperref[appendix-5f1a2c3f97ff]{Numerical reproducibility} --- Rounding, contraction, subnormal and optimization assumptions must match.
\item
  \hyperref[appendix-e979eb012782]{Flush-to-zero and denormals-are-zero} --- uses this mechanism in its explanation.
\item
  \hyperref[appendix-1823a14330cb]{Interval arithmetic} --- uses this mechanism in its explanation.
\end{itemize}

\subsection{Recall and next depth}\label{appendix-f11d1a521253-recall-and-next-depth}

Reconstruct the concrete example in the definition and identify its validity assumptions.

\textbf{Pending:} Extend this mechanism to later algorithms and target-specific behavior. See the M005 chapter for its current examples and validity assumptions.

\subsection{Source and navigation}\label{appendix-f11d1a521253-source-and-navigation}

\hyperref[01-foundations-i-ch001-when-arithmetic-meets-uncertainty]{M005 development}

\hyperref[01-foundations-i-ch001-a-stretching-ruler-for-real-valued-quantities]{M004 teaching chapter} · \hyperref[appendix-fd090d982a15-m004---parts-61-76]{Per-part definitions} · \href{../07-Vision-XR-and-Playable-Worlds/main.pdf\#nameddest=appendix-4ad47bdd269a}{Dependency Map (Vision, XR, and Playable Worlds, Reference notes: Dependency Map)}
